% chapter2banner.tex
% Standalone chapter-opening page: badge + title + photo + objectives + intro.
% This page is entirely separate from chapter2.tex's lesson content.

\khmerornatetitle{ជំពូកទី២៖ រលក}{២. រលក}
\khmerbannerphoto{img/waves}

% --- CHAPTER LIST (2-column grid, edit freely) ---
\begin{center}
\begin{tabularx}{\textwidth}{@{} >{\raggedright\arraybackslash}X >{\raggedright\arraybackslash}X @{}}
  \khmercircnum{1}~ គោលការណ៍តម្រួតនៃរលក \mbox{ និងរលកជញ្ច្រុំ}
  &
  \khmercircnum{2}~ អាំងទែផេរ៉ង់នៃរលក
\end{tabularx}
\end{center}

\vspace{0.15cm}

% --- Intro paragraphs ---
\raggedright
\setlength{\parindent}{1.2em}
នៅពេលដែលអ្នកទាំងអស់គ្នាបើកភ្នែកនោះ អ្នកទាំងអស់គ្នានឹងសង្កេតឃើញនូវរូបភាពជាច្រើនដែលមានដូចជាពន្លឺ និងអ្វីៗគ្រប់យ៉ាង។ មិនថាសម្លេងនិងរំញ័រដែលយើងសង្កេតឃើញនៅក្នងជីវិតប្រចាំថ្ងៃសុទ្ធតែជារលកដែលយើងពុំអាចមើលឃើញ។ ជាក់ស្តែង យើងបានសិក្សាពីរលកទាំងពីថ្នាក់ទី១១ មកម្ល៉េះ ដែលធ្វើឲ្យយើងស្គាល់ពីសមីការលក ដែលយើងនឹងធ្វើការសិក្សាបន្ថែមនៅថ្នាក់ទី១២ទាក់ទងនឹងលក្ខខណ្ឌផ្សេងៗដូចជា គោលការណ៍តម្រួតនៃរលក រលកជញ្ច្រុំ និងអាំងទែផេរ៉ង់នៃរលក ដែលយើងសង្កេតឃើញមាននៅគ្រប់រូបភាពក្នុងជីវិតប្រចាំថ្ងៃ។ សូមស្វាគមន៍មកកាន់ជំពូកទី២។

\clearpage